\section{Introduction}
Cyber-Physical Systems (CPS)~\cite{baheti2011cyber} integrate computational and physical capabilities and are now common in safety-critical applications like autonomous driving and medical devices~\cite{rajkumar2010cyber}. For such systems, any mistake of them may result in catastrophic consequences. Linear Temporal Logic (LTL) is a formal language used to define the complex rules for safe operation, such as requiring a timely response or forbidding unsafe conditions~\cite{pnueli1977temporal}. Therefore, the ability to verify that a CPS follows its LTL specification is essential for building trust and ensuring safety in these systems.

Although important, formally verifying LTL specifications for CPS is a significant challenge. The main difficulty arises from the need to handle the infinite-horizon nature of LTL properties~\cite{baier2008principles} together with the infinite state spaces and complex dynamics of physical systems. To address this, researchers have developed certificate-based methods, including the state-triplet method~\cite{wongpiromsarn2015automata} and closure certificates~\cite{murali2024closure}. The state-triplet method provides an automata-theoretic approach for verifying LTL specifications of dynamical systems. This technique constructs barrier certificates that prevent transitions to accepting states in the automaton representing the negated specification. This approach can be conservative as it requires blocking all potential paths to accepting states individually. The method's strength lies in its direct connection to automata structure, but it may need numerous certificates for complex specifications. Closure certificates offer a generalized framework by operating on pairs of system states through transition invariants. This approach establishes certificates that enforce decrease conditions along trajectories, ensuring accepting states are visited only finitely often. The method reduces conservatism by allowing more flexible certificate structures and naturally handles liveness specifications. While these methods are significant, they each have trade-offs. The state-triplet method can be conservative and closure certificates have low synthesis efficiency because of the large search space they must explore. These limitations show a clear need for a new framework that can reduce conservatism without a high computational cost.

To address this gap, we propose a new framework to verify LTL specifications via proof certificates. Our approach shifts the focus of the LTL verification problem to the transitions of the Nondeterministic Büchi Automaton (NBA)~\cite{hahn2020reward}. In our framework, verification is achieved by finding two types of complementary certificates. Transition safety certificates prove that certain accepting transitions never occur, while transition persistence certificates ensure that other accepting transitions happen only a finite number of times. By combining them, we can formally guarantee that no accepting run exists. To find these certificates, we developed an automated synthesis method based on the Counterexample-Guided Inductive Synthesis (CEGIS) framework. This method supports both polynomial templates for efficiency and neural network templates for expressiveness. The main contributions of this paper are summarized as follows:
\begin{itemize}

  \item We propose a new LTL verification framework for discrete-time dynamical systems with continuous state spaces, which employs​ transition certificates to reduce the search space, thereby achieving higher computational efficiency compared with​ closure certificates.
  

  \item We introduce automated synthesis methods​ for transition certificates based on polynomial and neural network templates. These templates provide a balance between expressiveness and computational efficiency, enabling efficient certificate generation for diverse system classes, from linear to nonlinear dynamics.
  

  \item We demonstrate the effectiveness​ of our certificates through case studies, including the Kuramoto oscillator and a two-room temperature model. Experimental results show that transition certificates achieve significant speedups, highlighting their practicality for LTL specifications.
\end{itemize}

The organization of this paper is as follows: Section 2 covers preliminary knowledge. Section 3 introduces our proof certificates for LTL verification and compares with other methods. Section 4 presents synthesis methods. Section 5 demonstrates the effectiveness through case studies. Section 6 concludes the paper.


\textbf{Related work: } 
Verification of temporal logic properties via proof certificates has seen significant advancements. To validate discrete-time systems against LTL specifications, Wongpiromsarn et al.~\cite{wongpiromsarn2015automata} developed a state-triplet method that uses barrier certificates to block all simple paths from initial states to accepting states in the automaton, requiring enumeration of state triplets $(q_m, q_m', q_m'')$ along these paths. Murali et al.~\cite{murali2024closure} introduced closure certificates for persistence verification, which establish disjunctively well-founded transition invariants \cite{podelski2004transition} for dynamical systems, and further extended this notion to control closure certificates for $\omega$-regular specifications~\cite{murali2025control}. Additionally, substantial research has addressed the verification of specific temporal behaviors such as reachability, safety, and avoidance~\cite{chakarov2013probabilistic,chatterjee2016algorithmic,chatterjee2024sound,chatterjee2024equivalence}. Some adopt neural networks as templates~\cite{wang2025verifying, nadali2024neural, zhao2020synthesizing} to synthesize certificates for verifying temporal properties, including the simultaneous representation of ranking functions and inductive invariants~\cite{giacobbe2025let} and the verification of memoryful neural agents against LTL specifications~\cite{hosseini2025ltl}.


The abstraction-based approaches~\cite{vasile2016control, leahy2019control, kress2009temporal, coogan2015traffic, fainekos2009temporal, ding2014optimal, wongpiromsarn2009receding} lift the CPS model from continuous domain to discrete domain by computing a finite abstraction, e.g., finite transition systems and Markov decision processes (MDPs). We note that abstraction-based techniques have been used in the verification and synthesis of continuous-space dynamical systems against LTL properties such as the results in~\cite{henzinger1997hytech, khaled2021omegathreads, lahijanian2011temporal, rungger2016scots}. These results rely on building a finite state abstraction and then making use of model checking and synthesis techniques~\cite{baier2008principles, tabuada2009verification}. In general, the abstraction-based approaches are computationally demanding and suffer from the curse of dimensionality. Abstraction-free methods, such as proof certificates, can alleviate the computational burden to some extent.


%In the context of our discussion, we assume that the NBA for an LTL specification $\phi$ has already been constructed. The translation from an LTL formula to an NBA involves an exponential complexity with respect to the size of the formula~\cite{vardi2005automata}. The complexity of constructing the complement of an NBA is known to be EXPTIME~\cite{complementnba}. Converting LTL specifications into Transition-based Generalized Büchi Automata (TGBA) typically yields smaller automata compared to Labeled Generalized Büchi Automata (LGBA)~\cite{couvreur1999fly}. The degeneralization process~\cite{giannakopoulou2002states} can convert a TGBA with multiple acceptance conditions into a Transition-based Büchi Automaton (TBA) featuring a single acceptance condition. Various tools exist for generating different automaton types from LTL formulas, such as Owl~\cite{owl} and SPOT~\cite{spot}.
